\subsection{Интерпретация и ограничения}

Полученные результаты необходимо интерпретировать в логике исследования признаков, а не в логике соревнования моделей. Цель эксперимента состояла не в подборе архитектуры с максимальной точностью, а в проверке того, содержат ли FinBERT-generated features самостоятельный сигнал и дают ли они добавочную предсказательную силу относительно рыночных признаков. Такой подход согласуется с литературой по feature selection и feature contribution: новая группа признаков имеет смысл только тогда, когда её добавление улучшает качество прогноза при сопоставимом протоколе оценки \cite{guyon2003feature}.

Формально добавочную силу текстовых признаков можно определить через изменение выбранной метрики качества:

\[
\Delta_{text}^{M} = M(X_{price}, X_{text}) - M(X_{price}),
\]

где $M$ --- метрика качества, $X_{price}$ --- рыночные признаки, а $X_{text}$ --- текстовые признаки, полученные с помощью FinBERT. При такой записи положительное значение $\Delta_{text}^{M}$ означает, что добавление текстовых признаков улучшает результат относительно модели на рыночных признаках. Однако для устойчивого вывода этого недостаточно: желательно, чтобы объединённая модель также превосходила majority baseline и отдельную модель на текстовых признаках.

В RUN 03 для accuracy и balanced accuracy добавление FinBERT-признаков к рыночным признакам действительно даёт положительную разницу относительно \texttt{prices-only}:

\[
\Delta_{text}^{Accuracy} = 0.486 - 0.410 = 0.076,
\]

\[
\Delta_{text}^{Balanced\ Accuracy} = 0.474 - 0.400 = 0.074.
\]

Это означает, что текстовые признаки улучшают слабую модель, использующую только рыночные признаки. В этом смысле FinBERT embeddings оказываются не полностью случайной добавкой: при объединении с ценовыми признаками они повышают несколько метрик относительно \texttt{prices-only}. Такой результат согласуется с общей идеей textual finance, согласно которой тексты могут содержать информацию о настроениях, ожиданиях и интерпретации событий \cite{tetlock2007media,loughran2011liability,kearney2014textual,nassirtoussi2014text}.

Однако итоговый вывод должен быть осторожным. Во-первых, объединённая модель \texttt{prices + FinBERT} не превосходит majority baseline по accuracy. Во-вторых, ROC-AUC объединённой модели остаётся ниже 0.5:

\[
ROC\text{-}AUC_{combined} = 0.427 < 0.5.
\]

Это означает, что модель не показывает устойчивой способности ранжировать положительный и отрицательный классы. В-третьих, объединённая модель уступает \texttt{FinBERT-only}. Разность между объединённой моделью и текстовой моделью по accuracy отрицательна:

\[
\Delta_{combined-text}^{Accuracy} = 0.486 - 0.550 = -0.064.
\]

Следовательно, в текущем протоколе не подтверждается комплементарность рыночных и текстовых признаков. Текстовые embeddings демонстрируют слабый самостоятельный сигнал, но их объединение с рыночными признаками не приводит к устойчивому улучшению качества. Такой результат не противоречит литературе: прикладные работы показывают, что эффект текстовых данных зависит от источника сообщений, горизонта прогноза, способа агрегации текста и выбранного рынка \cite{hiew2019bert,zou2022prebit,jiang2023finbert,albada2025bertweet}.

Отдельно следует интерпретировать результат \texttt{FinBERT-only}. Эта модель показывает accuracy 0.550 против 0.532 у majority baseline. Разность составляет

\[
\Delta_{FinBERT-majority}^{Accuracy} = 0.550 - 0.532 = 0.018.
\]

С одной стороны, положительная разность указывает на наличие слабого самостоятельного текстового сигнала. С другой стороны, величина прироста мала, а ROC-AUC равен 0.512, то есть близок к случайному уровню. Поэтому корректная формулировка состоит не в том, что FinBERT-признаки «предсказывают рынок», а в том, что в выбранном эксперименте они дают слабый сигнал, недостаточный для сильного вывода о прогнозной устойчивости.

Результат LSTM также следует рассматривать через призму проверки признаков. LSTM использовалась как модель, способная учитывать последовательную структуру временного окна \cite{hochreiter1997lstm}. Если бы объединённые признаки содержали выраженный последовательный сигнал, можно было бы ожидать улучшения относительно логистической регрессии. Однако в RUN 03 LSTM на \texttt{prices + FinBERT} не превосходит простую модель. Это означает, что усложнение архитектуры не компенсирует ограниченную информативность признаков в текущей постановке. Аналогичный вывод встречается и в современных работах по LLM/FinBERT-признакам: более сложная языковая или нейросетевая модель не всегда превосходит простой baseline \cite{shobayo2024sentiment,vuong2025llm}.

Одно из возможных объяснений результата связано с размерностью признакового пространства. В объединённой модели используется

\[
d = 11 + 768 = 779
\]

признаков, тогда как число наблюдений в обучающей выборке равно $n_{train}=1104$. Отношение числа наблюдений к размерности признаков составляет

\[
\frac{n_{train}}{d} = \frac{1104}{779} \approx 1.42.
\]

Такое соотношение означает, что модель работает в относительно высокоразмерном пространстве при ограниченном числе наблюдений. Это может повышать риск переобучения и снижать устойчивость результата, особенно если embedding-признаки содержат шум или нерелевантные компоненты. Поэтому отсутствие устойчивого выигрыша объединённой модели может быть связано не только с бесполезностью текстов, но и с выбранным способом представления и агрегации текстовой информации.

Основные ограничения эксперимента можно разделить на несколько групп.

Во-первых, эксперимент проводится только на одном активе --- AAPL. Поэтому результат нельзя обобщать на весь фондовый рынок или на другие классы активов. В литературе показано, что вклад текстовых данных может различаться между акциями, индексами, криптовалютами и рынками с разной информационной структурой \cite{bollen2011twitter,mao2011sentiment,zou2022prebit}.

Во-вторых, используется один горизонт прогноза --- следующий торговый день. Текстовые признаки могут быть более полезны на других горизонтах: например, для внутридневной реакции на новости, для нескольких дней после события или для экстремальных движений. В работе PreBit, например, таргет связан с экстремальными движениями Bitcoin, а не с обычным next-day direction \cite{zou2022prebit}. Поэтому слабый результат на горизонте одного дня не исключает полезности LLM-generated features в другой постановке.

В-третьих, FinBERT используется как готовый генератор embedding-признаков без fine-tuning под конкретную задачу. Это соответствует цели работы, поскольку исследуется именно использование LLM/FinBERT-generated features как входных переменных. Однако в дальнейших экспериментах можно проверить fine-tuning, sentiment-score, topic features, emotion features или агрегацию нескольких текстовых представлений. Современные работы показывают, что разные способы извлечения признаков из текста могут давать разные результаты \cite{araci2019finbert,yang2020finbert,albada2025bertweet,vuong2025llm}.

В-четвёртых, качество текстового сигнала зависит от временного выравнивания. В данной работе используется консервативное правило: тексты для прогноза дня $t$ берутся не позднее $t-1$. Это снижает риск утечки, но может одновременно удалять часть релевантного новостного сигнала, опубликованного в день принятия решения. Более точная постановка потребовала бы внутридневных timestamp и разделения текстов на опубликованные до и после закрытия рынка.

В-пятых, результаты зависят от выбранного протокола оценки. В работе используется хронологическое разбиение без пересечения дат, что методологически строже случайного split. Однако для более устойчивой проверки можно использовать rolling-window validation. В такой схеме модель обучается на окне прошлых данных и тестируется на следующем периоде:

\[
Train_k = [t_0, t_k], \qquad Test_k = (t_k, t_{k+h}],
\]

где $h$ --- длина тестового окна. Такой подход позволил бы проверить, сохраняется ли эффект FinBERT-признаков в разных рыночных режимах.

Таким образом, результаты RUN 03 следует интерпретировать как осторожный отрицательный или слабоположительный вывод. FinBERT-признаки дают слабый самостоятельный сигнал и улучшают модель относительно \texttt{prices-only}, но не демонстрируют устойчивой добавочной предсказательной силы в более строгом сравнении с majority baseline и \texttt{FinBERT-only}. Поэтому основной вывод состоит не в том, что текстовые признаки бесполезны, а в том, что в выбранном MVP-протоколе их вклад оказался ограниченным и требует дальнейшей проверки на других активах, горизонтах и способах извлечения признаков.
